\chapter*{General Conclusion}
\addcontentsline{toc}{chapter}{General Conclusion}
\markboth{General Conclusion}{}

This research internship at the Maersk Mc-Kinney Moller Institute
produced a complete, reproducible pipeline for generating visual
counterfactual explanations on chest radiographs: RSNA data ingestion, a
Stable Diffusion~1.5 generator adapted to the radiograph domain through
LoRA at 0.78\% of its parameters, a gradient-guided DDIM sampler with a
hard segmentation-gated background lock, two directional PPO policies
controlling that sampler step by step, and an evaluation suite following
the protocol of Singla~\emph{et al.} --- all executable as batch jobs on
an offline HPC cluster.

Measured against the guidance classifier, the system reaches high
counterfactual validity (mean CV $0.975$) with good identity
preservation, and the background lock demonstrably confines edits to the
lung field.

The internship's most important result, however, is the one that
undermines that headline number. An independent held-out validation ---
re-scoring the generated counterfactuals with classifiers that never
participated in the guidance loop --- showed that counterfactuals
accepted 99.5\% of the time by the guidance model were accepted at most
8\% of the time by an independent one, while an agreement control
confirmed the validators understood the task on unedited images. The
system had, in significant part, learned to fool its own critic rather
than to depict or remove disease.

That finding reframed the second half of the internship. Ensemble
guidance, a Grad-CAM-derived constraint on the action space, and an
annealed sparsity reward were designed, implemented and enabled
specifically to attack the adversarial failure mode. Their validation
was not completed within the internship period, and this report is
explicit that their effectiveness remains an open question rather than a
demonstrated improvement.

A further, later phase of work (Chapter~\ref{ch:integrity}) turned the same
falsification discipline on the evaluation apparatus itself, auditing the
control-arm and held-out-transfer machinery built to answer that open
question. The audit found, and closed, three uncontrolled confounds in the
measurement pipeline and two silent bugs that would each have invalidated a
comparison built on top of them --- including a routing bug that would have
made every control arm's output identical, and a denominator bug that had
inflated every previously published transfer ratio. It also found, and
corrected, a set of FID values in the project's own recorded results that
had been altered by hand. Of the two countermeasures, ensemble guidance is
now shown, across three independent seeds, to modestly but reliably improve
lesion targeting and, at a larger ensemble size, to substantially improve
held-out transfer; the Grad-CAM spatial constraint is shown instead to hurt
lesion targeting rather than help it. A full held-out re-validation of the
surviving countermeasure was the item this chapter did not close, and is
taken up in Chapter~\ref{ch:closing}.

Chapter~\ref{ch:closing} closes it, partially. Both directional policies
were retrained from scratch under ensemble guidance --- not merely
evaluated through it, as the no-policy arms above were --- and re-scored
against the leak-free validator set. Mean held-out transfer improves on
both directions, \pton{} very nearly doubling; the trained-vs-random gap
from Table~\ref{tab:lhr_arms} is confirmed statistically significant
($p=0.012$) at a small effect size. But the single most independent
validator, \texttt{nih}, does not improve on \ntop{}, remaining in the
adversarial range. Ensemble guidance is accordingly reported as a real,
partial countermeasure, not a solved problem --- the same falsification
discipline this report has applied to every favourable number applied to
its own fix, and did not let it pass uncontested.

Technically, the project provided sustained exposure to latent diffusion
models, parameter-efficient fine-tuning, reinforcement learning with
continuous action spaces and shaped rewards, and the practical discipline
of research computing on a shared, non-interactive, walltime-bounded
cluster --- where an experiment must be fully specified in advance,
instrumented well enough to be diagnosed from a log file, and cheap
enough to fail fast.

As additional work beyond the assigned scope, the pipeline was extended to
a second, deliberately dissimilar imaging domain --- mammography, on
CBIS-DDSM (Chapter~\ref{ch:breast}) --- to test whether this lesson was
specific to RSNA or would recur elsewhere. It recurred: a mean
Counterfactual Validity of $1.000$ against the mammography guidance
ensemble was accompanied by a significant density shift with no
corresponding change in lesion size, the same diffuse-edit signature the
RSNA held-out validation had exposed, this time caught by a purpose-built
diagnostic metric rather than by an independent classifier, since training
one for that domain remains future work. This is reported as additional,
less-validated evidence alongside the RSNA result, not as a second
demonstrated success, but it strengthens rather than weakens the report's
central claim: a system's own validity number is not to be trusted without
an independent check built to falsify it.

The methodological lesson runs deeper than any of those techniques. The
most dangerous metric in a machine learning project is the one the system
was optimised against. Counterfactual validity computed with the guidance
classifier was not a weak measurement --- it was a circular one, and it
reported near-perfect success for a system that was substantially
failing. Every subsequent design decision of value in this project
descended from building the control capable of exposing that, and from
being willing to report what it found. For an internship whose purpose
was to bridge academic training and research practice, learning to
distrust a favourable number was worth more than obtaining one.

\clearpage
